\section{Nonempty tangent neighborhoods}
\label{sec:nonempty87}
A prescribed remainder allowance need not admit any legal perturbation.
The following quantitative realization makes the nonvacuous range
explicit.  It is a sufficient allowance, not a minimal one.  The
underlying positive-semidefinite tangent condition is standard
semidefinite variational geometry; see \cite{BCS87}.  What is needed
here is one common normalization and a finite remainder bound for the
whole measurement tuple.

\begin{proposition}[A sufficient allowance]\label{prop:nonempty87}
Let $E$ be an ordered $k$-outcome measurement and let $H$ be a
Hermitian zero-sum tuple with $Q_jH_jQ_j\succeq0$ for every $j$.
Choose $a>0$ and $\lambda>0$ such that
\[
 \|H_j\|_{\mathrm{op}}\le a,\qquad
 E_j\succeq\lambda P_j\quad(1\le j\le k),
 \qquad P_j=\operatorname{supp}E_j.
\]
Set
\begin{equation}\label{eq:nonemptyconstants87}
 c=1+2a^2/\lambda,\qquad
 s_* =\min\{1,\lambda/(2a),1/a\},\qquad
 \Lambda_* =ck(1+2k).
\end{equation}
Then, for every $0\le s\le s_*$,
\begin{equation}\label{eq:nonemptycurve87}
 F_j^{\mathrm c}(s)=\frac{E_j+sH_j+cs^2I_d}{1+kcs^2}
\end{equation}
is a legal measurement and satisfies
\begin{equation}\label{eq:nonemptybound87}
 \|F^{\mathrm c}(s)-E-sH\|_{\Sigma}\le\Lambda_*s^2.
\end{equation}
Thus every allowance $\Lambda\ge\Lambda_*$ gives a nonempty
neighborhood at every $s\in(0,s_*]$.  For arbitrary real input,
\eqref{eq:nonemptycurve87} is a rational function of $s$ with matrix
coefficients.  For Gaussian-rational $E,H$ one may choose rational
$a,\lambda,c,s_*,\Lambda_*$, giving Gaussian-rational coefficients.
\end{proposition}
\begin{proof}
For $P_j\ne0$, the support block of the numerator in
\eqref{eq:nonemptycurve87} is at least $\lambda P_j/2$.
Its cross block is $sP_jH_jQ_j$.  The Schur complement on the
missing support is at least
\[
 sQ_jH_jQ_j+s^2(c-2a^2/\lambda)Q_j\succeq0.
\]
If $P_j=0$, the hypothesis says $H_j\succeq0$, so the numerator
is positive directly.  This also covers zero outcomes without
assigning them a smallest positive eigenvalue.  The common denominator
is positive for all real $s$, and summing the numerators gives exactly
$(1+kcs^2)I_d$.

Direct subtraction, with no series expansion, gives
\[
 F_j^{\mathrm c}(s)-E_j-sH_j
 =\frac{cs^2(I_d-kE_j-ksH_j)}{1+kcs^2}.
\]
Every effect has norm at most one, and $sa\le1$ on the given
interval.  The component remainder is therefore at most
$c(1+2k)s^2$.  Summing proves \eqref{eq:nonemptybound87}.
There is a positive common support floor since there are finitely
many effects and at least one is nonzero.  Rational minorants and
majorants give the last assertion.
\end{proof}

\begin{proposition}[Exact nonemptiness certificates]\label{prop:nonemptyexact87}
On Gaussian-rational represented input, a sufficient certificate
\eqref{eq:nonemptyconstants87} can be constructed and verified in
polynomial bit complexity.  It certifies the whole explicit curve
\eqref{eq:nonemptycurve87}, not only sampled values.  Its interval
$s_*$ is an interval of legality and remainder control; it is not a
computed small interval for the discrimination laws.
\end{proposition}
\begin{proof}
Use the rational support projectors of
Proposition~\ref{prop:supportexact84}.  The maximum absolute row sum,
with $|\Re z|+|\Im z|$ in place of $|z|$, gives a rational upper
bound $a\ge1$ for every Hermitian $H_j$.
Starting at $\lambda=1$, halve until the exact positive-semidefinite
tests $E_j-\lambda P_j\succeq0$ all hold.  Nonzero eigenvalues of a
rational positive matrix have a lower bound $2^{-\operatorname{poly}(d,L)}$
in terms of its represented bit length $L$: clear denominators and
use the nonzero coefficient of lowest degree in the characteristic
polynomial, together with the row-sum bound on the other eigenvalues.
Thus only polynomially many halvings are required.  All remaining
operations in \eqref{eq:nonemptyconstants87} are rational arithmetic.
The displayed PSD predicates and algebraic proof above verify the
certificate on the continuum.  A resource-limited implementation
must reject an exhausted cap, rather than return a partial certificate.
No claim that $\Lambda<\Lambda_*$ is infeasible follows from this
sufficient construction.
\end{proof}
